上一节把跨组协调显式放在购买服务中。分布式数据库可以把同类责任放入事务层，让一个SQL事务修改多个数据分片。整个结构同时回答几个问题：数据键放在哪里，每个分片怎样复制，多个分片怎样共同决定，事务间怎样并发，以及读取哪个时间的版本。

\subsection{Spanner怎样组合数据与复制组}

Spanner面向跨数据中心保存数据，并向应用提供有模式的表、查询和通用事务。表中的行由有序主键定位，底层把行列数据组织为带版本的键值映射。相关行的键前缀与放置规则一起决定数据局部性。购买绑定、购买记录和原C/T若常被一起访问，把它们放近，可以减少一次操作跨越的组数~\cite{spanner2012}。

存储服务器spanserver承载tablet。这里的tablet是保存一组键值映射的容器，在本地包含类似B树的文件结构和预写日志，底层字节由分布式文件系统保存。每个tablet上有一个Paxos状态机，跨机器的对应副本共同组成Paxos组，复制同一tablet的状态。一个机器可以承载许多tablet，每个组则独立处理自己的复制。

组的leader组织需要协调的写入和事务读。leader上的锁表记录键范围的锁状态，事务管理器参与跨组提交，事务管理器状态也通过所属Paxos组保存。leader退出后，后继者从复制事实恢复协调责任。副本中的数据、组内复制记录、锁及跨组事务阶段，各自有清楚的保存或重建位置。

\CNineteenFigure{ch19-13}{Spanner在各Paxos复制组之上组织锁表与跨组事务}{fig:19-spanner}

在键值映射上，Spanner还组织directory，即共享前缀的一组连续键。directory表达应用数据的局部性与放置要求，一个Paxos组可以承载多个directory。因此，一个tablet可以容纳多个不连续的键区间。较大的directory再分成fragment，实际移动可以按fragment进行。这里的directory是数据放置抽象，本书定位逻辑分区承载组的路由目录则负责另一类元数据。

\CNineteenFigure{ch19-28-spanner-layout}{键前缀形成放置关系，tablet容纳directory的fragment，Paxos组复制该容器}{fig:19-spanner-layout}

相关数据都落在一个组时，事务可以借助该组的锁表和Paxos完成。跨多个组时，各组leader执行两阶段提交，其中一个参与组承担协调组。数据放置使常见事务留在组内，通用跨组路径则处理确实同时依赖多个分片的业务。

数据移动先登记迁移，后台搬运主要内容，最后用事务移动少量剩余数据并原子更新相关元数据。把大规模字节传输与最终交接分开，可以缩短阻断在线请求的阶段。在线移动期间，事务继续修改数据，最后的交接必须包含搬运期间的更新；第19.2节的冻结图像则先停止旧前缀变化，再写齐块并发布清单。两条路径据各自的锁、事务与资格规则确定可以切换的位置~\cite{spanner2012}。

\subsection{复制、prepare与共同提交时间}

读写事务用两阶段锁保护读写对象，写入先在客户端缓冲。提交开始后，非协调参与组取得写锁，选择prepare时间戳，通过本组Paxos保存prepare记录，再把时间戳告知协调组。协调组取得自己的写锁，收齐必要准备以后，为整个事务选择提交时间$s$，并通过本组Paxos保存commit记录。

$s$满足参与者prepare时间戳、本组时间戳单调推进和时间接口给出的下界。所有组按同一个$s$形成多版本写入，数据库便可以在一个共同时间上解释事务结果。参与组知道共同决定以后，再通过本组复制保存结果，按$s$应用并释放锁。

这里有两种不同的日志事实：参与者保存prepare，以及事务最终提交购买或库存写入。前者已经进入持久复制，后者还可能等待共同决定。读取复制日志位置时，先解释该位置所保存的阶段，才能判断业务数据是否已经提交。

\subsection{TrueTime怎样表达时间不确定性}

事务的串行顺序还要与现实中已经确定的先后相容。Alice取得事务T1的成功结果以后，再发起T2，外部一致性要求T1排在T2之前。机器时钟存在误差，仅把本地时钟数值作为提交时间，可能让后开始的事务拿到更早数值。Spanner把误差区间直接交给提交协议处理。

TrueTime的\code{TT.now()}返回区间$[\mathrm{earliest},\mathrm{latest}]$，保证调用时刻的真实时间位于其中。区间宽度给出当前不确定性，应用得到的是可用于判断先后的范围。

\code{TT.after(s)}则表示可以确定真实时间已经越过$s$，其判断依赖当前区间的earliest已经大于$s$~\cite{spanner2012}。

底层由时间参考、机器时钟同步、保守漂移界和异常时钟排除共同维护该保证。论文使用GPS与原子钟等不同参考，机器定期向多个时间主机校准，在两次校准之间扩大误差界。参考不可达、机器过载或网络延迟会使区间变宽，提交协议据此增加等待。

\CNineteenFigure{ch19-29-tt-interval}{提交时间取在当前区间及prepare下界之后，等待新的earliest越过它}{fig:19-tt-interval}

图~\ref{fig:19-tt-interval} 中，设当前区间为$[100,108]$，参与者prepare下界为105，协调者可以选择$s=109$，前提是还满足本组的单调要求。随后区间变为$[110,118]$，earliest已经大于109，\code{TT.after(109)}成立。

协调者选$s$时，采用收到提交请求后取得的\code{TT.now().latest}作为下界之一。保存commit决定以后，在允许提交可见以前等到\code{TT.after(s)}成立，这个阶段称为commit wait，即提交等待。它可与Paxos通信部分重叠，然后协调者向客户端及其他参与组传播完成结果。

\CNineteenFigure{ch19-14}{选择时间上界并等待其过去，使提交次序符合现实中已经分开的事务先后}{fig:19-external-consistency}

设T1已经向客户端完成，客户端随后才开始T2。T1完成以前，真实时间已超过$s_1$。T2的提交请求又晚于T2开始，选出的$s_2$不早于请求到达后取得区间的上界。于是
\[
s_1<t_{\text{T1完成}}<t_{\text{T2开始}}
\le t_{\text{T2提交请求到达}}\le s_2\text{，}
\]
得到$s_1<s_2$。对时间上重叠的事务，锁与原子提交决定合法串行顺序。若区间错误地排除真实时间，上述不等式便失去依据，时间参考和误差维护属于该结构的成立条件。

\subsection{固定时间读取与safe time}

多版本数据使读取者可以请求某个时间$t$的数据库状态。只读事务先选定读取时间，再从足够新的副本取得该时点版本，无须取得读写事务使用的锁。一个副本是否足够新，还取决于它有没有遗漏可能在$t$或更早提交的事务。

Spanner为副本记录safe time，即当前可安全服务的最大时间。它同时受两个进度限制：
\[
t_{\mathrm{safe}}=\min(t_{\mathrm{Paxos}},t_{\mathrm{TM}})\text{。}
\]
$t_{\mathrm{Paxos}}$来自按序应用的复制写入进度。$t_{\mathrm{TM}}$来自事务管理器保存的未决prepare：某事务在$p$准备，其最终提交时间可能从$p$开始，因此该事务决定以前，可服务时间必须停在$p$之前。没有未决prepare时，这一项不额外限制~\cite{spanner2012}。

\CNineteenFigure{ch19-30-safe-time}{复制应用已到130，未决prepare仍可能在110提交，安全读取先停在110之前}{fig:19-safe-time}

图中副本已经应用到130，仍有prepare下界为110的事务。它可以服务时间100的读取，时间115的读取则等待共同决定和相应应用。另一条更晚写入已应用，并没有说明这项未决事务最终会commit还是abort。决定到达后，事务管理器解除对应限制，副本再按新的安全上界服务。

全库对账可以固定一个读取时间，逐组取得一致观察。查询遇服务器故障时携带同一读取时间与当前位置，在另一个满足safe time条件的副本上继续。作业保存这个时间及输入范围，在线新付款继续发生，固定输入的身份保持明确。

\subsection{CockroachDB的Range、意图与事务记录}

CockroachDB把SQL数据编码为有序键值，将键空间切成Range，每个Range由Raft组复制。SQL网关把事务操作送往对应Range，分布式事务层连接这些Range上的读写。Range的副本保存数据，leaseholder承担相应的当前读写协调职责，事务记录则关联跨Range的未完成写入。本节按v25.4.0及同版架构文档说明~\cite{crdb254-transactions}。

事务写入先形成write intent，即带所属事务信息的暂定值。其他操作遇到它，通过事务状态判断等待、冲突处理或解析意图。事务记录保存PENDING、STAGING、COMMITTED或ABORTED，多个意图指向同一记录。各Range的Raft负责保存意图与记录，事务协议据这些事实解释整体结果。

Parallel Commits把事务记录推进到STAGING，并检查登记的在途写是否已经成功复制。满足协议的全部条件时，事务隐式提交，可以回复成功。随后异步把记录推进到COMMITTED，把意图解析为普通多版本值。STAGING单独出现时，还要核对它引用的写入状态~\cite{crdb254-transactions}。

\CNineteenFigure{ch19-15}{事务记录与它引用的在途写共同说明提交状态，异步清理随后完成}{fig:19-crdb}

CockroachDB用混合逻辑时钟HLC给版本和事务分配时间戳。HLC包含接近墙钟的物理分量，以及区分同一物理分量下事件的逻辑分量。节点发消息时携带HLC，接收方推进自己的时钟，使之后的事件排在已观察事件之后。它把机器时间与因果传播组合为可排序时间戳。

\CNineteenFigure{ch19-31-hlc}{消息携带HLC，接收方推进本地时间戳，真实时间误差另外由时钟同步和偏移规则约束}{fig:19-hlc}

HLC时间戳缓存、锁与读集刷新共同处理并发。写入时间若被推进，SERIALIZABLE事务检查已经读过的范围能否在新时间保持有效，刷新失败则要求重试。该版本也提供READ COMMITTED等隔离选择，本书依赖可串行化的购买事务应明确使用SERIALIZABLE。

HLC本身没有返回TrueTime式的真实时间区间。产品通过时钟同步、最大偏移配置和节点自检约束现实时间偏差。官方文档分别说明SERIALIZABLE的串行化保证，以及偏差越界可能影响因果相关事务间单键线性化的情况。读取一致性还要按所选隔离级别和读路径判断~\cite{crdb254-transactions}。

每个Range的closed timestamp表示今后不再引入不晚于该时间的新写。follower应用到随此承诺关联的Raft位置后，可以服务该时间及更早的历史读取。当前读与允许陈旧的follower read据不同条件选择副本，历史读取的低延迟来自已经闭合的过去区间。

\subsection{TiDB事务怎样跨越TiKV Region}

TiDB负责SQL与事务上下文，TiKV负责有序KV及Region内复制，PD提供时间戳、位置与调度等服务。第17.5节已经给出Region和Peer的关系：一个Region覆盖键范围，多个Peer组成它的Raft组。现在一个事务的键跨过多个Region，上层协议把它们关联为共同提交~\cite{tikv-75-storage,tidb-75-pd}。

事务从PD取得\code{start\_ts}，以它选择快照版本，修改先积累在TiDB的事务上下文。基础两阶段提交选一个待写键作为事务primary key，按Region分组发送prewrite。TiKV检查冲突，保存待写数据和锁，锁里包含primary关联~\cite{tidb75-commit}。

此处primary key表示本次提交协议的主记录键。SQL表主键标识一行，Region leader组织一个范围的复制，事务primary关联本次跨Region决定，三者各有对象。另一个事务可以选不同键作为自己的primary。

\CNineteenFigure{ch19-16}{基础两阶段提交跨多个Region，事务primary关联各键的最终结果}{fig:19-tidb}

全部prewrite成功以后，事务取得\code{commit\_ts}并提交primary。primary成功后可以回复客户端，其余secondary的提交随后继续。读取者遇到未解析锁，按primary的状态判断：已提交就按原\code{commit\_ts}推进secondary，明确回滚则按原事务清除相应准备。事务未终结时，按产品锁与事务状态协议继续恢复。

\CNineteenFigure{ch19-32-mvcc}{同一事务的锁关联primary；primary已提交以后，secondary锁的解析继续形成同一版本}{fig:19-tidb-mvcc}

TiDB节点退出以后，已经落入TiKV的锁、数据与提交记录承接原事务关系。一个Region保存了prewrite，只能说明本地相关键已经准备，另一Region的冲突仍可能使整个事务回滚。primary已经提交而secondary尚有锁时，锁的存在表示清理未完成，读取者从原提交事实推进。

TiDB v7.5还提供Async Commit与单阶段提交优化，配置和事务条件决定实际路径。前者异步推进相应提交工作，后者在适用的单Region条件下减少阶段。新建集群默认启用这两个开关，升级集群则要检查沿用的配置，事务再按条件选择路径。分析故障前先确认实际模式，再按该协议解释记录~\cite{tidb75-commit}。

新建v7.5集群默认采用悲观事务，SQL也可以显式选择乐观或悲观模式。悲观模式在相关写入或\code{SELECT FOR UPDATE}时取得锁，普通快照读取保持版本语义。事务模式说明何时取得锁，隔离级别说明允许怎样的并发历史，判断跨行不变量还要核对共同读写对象。

TiDB为兼容MySQL将快照隔离标为REPEATABLE-READ，快照隔离允许写偏斜。设A、B仓库各有一件应急库存，规则要求合计至少留一件。两个事务读同一快照，都看到合计两件，一个只预留A，另一个只预留B。它们修改不同的行，逐行写冲突检查可能都通过，合计库存变成零~\cite{tikv-75-transactions}。

\CNineteenFigure{ch19-33-write-skew}{两个事务写不同的行，快照隔离下共同读取的总量约束可能被同时破坏}{fig:19-write-skew}

应用可以把共同约束集中到一个保护记录，或显式锁定共同判断依赖的行，使两次并发转移在相同对象上协调。TiDB范围查询的锁又有具体产品边界，不能直接套用InnoDB间隙锁的行为。各Region复制与跨Region原子提交保存共同结果，并发规则负责决定这个结果是否满足业务读写条件。

\subsection{事务中的购买身份}

数据库能为纳入事务的数据提供原子性和相应隔离，应用决定共同保存哪些事实。恢复原购买需要完整$S$、原输入、实际P/C/T和原结果，并用约束或仲裁让同一$S$回到同一记录。两次独立事务各扣一次库存，在数据库看来可以是合法串行历史，业务身份说明它们其实来自同一购买的重试。

提交超时后，客户端按原$S$查询或重入仲裁。最终已提交就读取原记录和结果，最终未提交则由合法尝试建立绑定。可能重新执行的事务代码只安排可放弃的读取、计算和待提交写入。实际付款从持久C/T取得资格，用原平台身份发送，数据库冲突重试便不会把一次外部效果当作可撤销的内存计算再次执行。
\FloatBarrier
